%-------------------------
% Resume in Latex (Improved Project Entry)
% Author : Jake Gutierrez (original) / Edited by OpenCode
% License : MIT
%------------------------

\documentclass[letterpaper,11pt]{article}

\usepackage{latexsym}
\usepackage[empty]{fullpage}
\usepackage{titlesec}
\usepackage{marvosym}
\usepackage[usenames,dvipsnames]{color}
\usepackage{verbatim}
\usepackage{enumitem}
\usepackage[hidelinks]{hyperref}
\usepackage{fancyhdr}
\usepackage[english]{babel}
\usepackage{tabularx}
\input{glyphtounicode}

\pagestyle{fancy}
\fancyhf{} % clear all header and footer fields
\fancyfoot{}
\renewcommand{\headrulewidth}{0pt}
\renewcommand{\footrulewidth}{0pt}

% Adjust margins
\addtolength{\oddsidemargin}{-0.5in}
\addtolength{\evensidemargin}{-0.5in}
\addtolength{\textwidth}{1in}
\addtolength{\topmargin}{-.5in}
\addtolength{\textheight}{1.0in}

\urlstyle{same}

\raggedbottom
\raggedright
\setlength{\tabcolsep}{0in}
\setlength{\footskip}{12pt}

% Sections formatting
\titleformat{\section}{
  \vspace{-4pt}\scshape\raggedright\large
}{}{0em}{}[\color{black}\titlerule \vspace{-5pt}]

% Ensure that generate pdf is machine readable/ATS parsable
\pdfgentounicode=1

%-------------------------
% Custom commands
\newcommand{\resumeItem}[1]{
  \item\small{
    {#1 \vspace{-2pt}}
  }
}

\newcommand{\resumeSubheading}[4]{
  \vspace{-2pt}\item
    \begin{tabular*}{0.97\textwidth}[t]{l@{\extracolsep{\fill}}r}
      \textbf{#1} & #2 \\
      \textit{\small#3} & \textit{\small #4} \\
    \end{tabular*}\vspace{-7pt}
}

\newcommand{\resumeSubSubheading}[2]{
    \item
    \begin{tabular*}{0.97\textwidth}{l@{\extracolsep{\fill}}r}
      \textit{\small#1} & \textit{\small #2} \\
    \end{tabular*}\vspace{-7pt}
}

\newcommand{\resumeProjectHeading}[2]{
    \item
    \begin{tabular*}{0.97\textwidth}{l@{\extracolsep{\fill}}r}
      \small#1 & #2 \\
    \end{tabular*}\vspace{-7pt}
}

\newcommand{\resumeSubItem}[1]{\resumeItem{#1}\vspace{-4pt}}

\renewcommand\labelitemii{$\vcenter{\hbox{\tiny$\bullet$}}$}

\newcommand{\resumeSubHeadingListStart}{\begin{itemize}[leftmargin=0.15in, label={} ]}
\newcommand{\resumeSubHeadingListEnd}{\end{itemize}}
\newcommand{\resumeItemListStart}{\begin{itemize}}
\newcommand{\resumeItemListEnd}{\end{itemize}\vspace{-5pt}}

%-------------------------------------------
%%%%%%%  RESUME STARTS HERE  %%%%%%%%%%%%%%%%%%%%%%%%%%%%

\begin{document}

%----------HEADING----------
\begin{center}
    {\Huge \textsc{MD SOHANUZZAMAN SHANTO}} \\
    \vspace{6pt}
    {\small +8801774786111 $|$ \href{mailto:md.sohanuzzaman@proton.me}{\underline{md.sohanuzzaman@proton.me}} $|$ \\
    \href{https://linkedin.com/in/sohanuzzaman3301}{\underline{linkedin.com/in/sohanuzzaman3301}} $|$ \\
    \href{https://github.com/sohanuzzaman3301}{\underline{github.com/sohanuzzaman3301}} $|$ \\
    \href{https://sohanuzzaman3301.github.io}{\underline{sohanuzzaman3301.github.io}}}
\end{center}

%-----------EDUCATION-----------
\section{Education}
\resumeSubHeadingListStart
\resumeSubheading
{Wuhan Institute of Technology}{Wuhan, China}
{Bachelor of Engineering in Computer Science and Technology}{Feb. 2022 -- Jan 2026}
\resumeItemListStart
\resumeItem{Grade: 84\% | Awards: First Class Scholarship for Academic Excellence (Ranked 1st in batch, 2025)}
\resumeItemListEnd
\resumeSubHeadingListEnd

%-----------CERTIFICATIONS-----------
\section{Certifications \& Licenses}
\resumeSubHeadingListStart

\resumeSubheading
{SOC Level 1 Certificate}{Mar. 2026}
{TryHackMe}{}
\resumeItemListStart
\resumeItem{Performed security monitoring, alert triage, incident investigation, and threat detection in simulated SOC environments.}
\resumeItem{Utilized Splunk, Elastic Stack (ELK), Wireshark, Zeek, Sysmon, and Wazuh for log analysis, network monitoring, and endpoint investigations.}
\resumeItem{\small\href{https://tryhackme-certificates.s3-eu-west-1.amazonaws.com/THM-NXIMKAYVCK.pdf}{TryHackMe certificate}}
\resumeItemListEnd

\resumeSubheading
{Google Cybersecurity Professional Certificate}{Sep. 2025}
{Google / Coursera}{}
\resumeItemListStart
\resumeItem{SIEM monitoring, incident response, network traffic analysis, Python automation, and NIST security frameworks.}
\resumeItem{\small\href{https://www.credly.com/badges/ed2ab68a-c056-4cd1-a293-82a4d2422a44/public_url}{Credly badge}}
\resumeItemListEnd

\resumeSubheading
{Google IT Support Professional Certificate}{Nov. 2025}
{Google / Coursera}{}
\resumeItemListStart
\resumeItem{System administration, Active Directory, networking, and IT infrastructure troubleshooting.}
\resumeItem{\small\href{https://www.credly.com/badges/98cc58ec-2e39-4863-938f-b5c895a6dcb8/public_url}{Credly badge}}
\resumeItemListEnd

\resumeSubHeadingListEnd

%-----------PROJECTS-----------
\section{Projects}
    \resumeSubHeadingListStart
      \resumeProjectHeading
          {Home Lab — SOC Detection \& Incident Response (Independent)}{2025 -- 2026}
          \resumeItemListStart
            \resumeItem{Built a production-like SOC home lab: deployed Splunk Enterprise (indexers, search head), Elastic Stack for historical storage, a Zeek+Suricata network sensor, and Windows Domain (AD) + instrumented endpoints (Sysmon, Winlogbeat). Collected: Windows Event Logs, Sysmon (EventID 1/3/7/11/12), Zeek conn/http/dns logs, and PCAPs for end-to-end visibility.}
            \resumeItem{Detection engineering: authored and validated actionable detection rules in Splunk SPL and Microsoft Sentinel KQL, and translated detections to Sigma for portability. Example Splunk SPL (SMB lateral movement): \texttt{index=wineventlog EventCode=4688 CommandLine="*\\psexec.exe*" | stats count by Computer,User,CommandLine}.}
            \resumeItem{Network triage workflow: extracted Zeek HTTP logs, used Wireshark to inspect PCAPs and identify beaconing (periodic HTTP POSTs). Applied Wireshark filters for focused analysis: \texttt{http.request.method == "POST" and tcp.len > 0} and used Zeek scripts to extract domains and user-agents for IOC creation.}
            \resumeItem{Endpoint hunt \& pivot: correlated Splunk process creation (Sysmon EventID 1) with network connections (Sysmon EventID 3) to map parent-child relationships, identified anomalous rundll32 and cmd invocations, and traced origin to phishing-delivered payload.}
            \resumeItem{ATT\&CK mapping \& remediation: mapped observed techniques to MITRE ATT\&CK (T1071.001 — Web Protocols, T1021.002 — SMB/Windows Admin Shares, T1547.001 — Service Persistence). Produced a remediation runbook with containment steps (host isolation, service disablement, credential reset), IOC blocklists (IPs/domains), and suggested firewall/IDS rules for prevention.}
            \resumeItem{Deliverables: produced a detection pack comprising 6 Splunk saved searches, 4 KQL analytics rules, 3 Sigma rules, Zeek log parsers, example YARA rules for identified malware, and a one-page incident report template mapping findings to ATT\&CK techniques for quick stakeholder briefing.}
            \resumeItem{Verification \& automation: replayed PCAPs with tcpreplay to validate detections, automated log ingestion via Winlogbeat/Beats -> ELK, and scripted recurring validation using Python to assert detection coverage against a test corpus of benign and malicious traces.}
          \resumeItemListEnd
      \resumeProjectHeading
          {Password Strength Meter $|$ \emph{Astro, JavaScript, CSS}}{2024}
          \resumeItemListStart
            \resumeItem{Client-side password strength estimator and leak checker that estimates entropy and likely crack times using a zxcvbn-inspired approach.}
            \resumeItem{Performs k-anonymity "Have I Been Pwned" leak checks entirely from the browser so plaintext passwords are never transmitted; includes CI for deployment to Cloudflare Pages.}
          \resumeItemListEnd
      \begin{itemize}[leftmargin=0.15in, label={}]
      \item \small GitHub: \href{https://github.com/Sohanuzzaman3301/password-strengh-meter}{https://github.com/Sohanuzzaman3301/password-strengh-meter}
      \item \small Live: \url{https://dd3b2174.password-strength-meter.pages.dev/}
      \end{itemize}
    \resumeSubHeadingListEnd

%-----------RESEARCH EXPERIENCE-----------
\section{Research Experience}
    \resumeSubHeadingListStart
      \resumeSubheading
          {Undergraduate Thesis: Explainable PE Malware Detection using Ensemble Learning}{Wuhan Institute of Technology}
          {EMBER2024 dataset (1M+ PE samples); ensemble stacking and SHAP-based explainability for malware classification}{2024 -- 2026}
          \resumeItemListStart
            \resumeItem{Developed an ensemble stacking pipeline for PE malware classification trained on the EMBER2024 dataset (1M+ samples).}
            \resumeItem{Applied SHAP to produce model-agnostic explanations and feature-level attributions to improve analyst interpretability and debugging.}
          \resumeItemListEnd
      \begin{itemize}[leftmargin=0.15in, label={}]
      \item \small GitHub: \href{https://github.com/Sohanuzzaman3301/pe-malware-ensemble/}{https://github.com/Sohanuzzaman3301/pe-malware-ensemble/}
      \end{itemize}
    \resumeSubHeadingListEnd

%-----------VIRTUAL EXPERIENCE & PROJECTS-----------
\section{Virtual Experience \& Projects}
\resumeSubHeadingListStart
\resumeSubheading
{Mastercard Cybersecurity Job Simulation | Forage}{June 2026}
{Remote}{}
\resumeItemListStart
\resumeItem{Simulated the role of a Mastercard Security Analyst to design and launch realistic, company-wide phishing simulation drills.}
\resumeItem{Analyzed employee vulnerability data to identify high-risk departments and security gaps across the enterprise.}
\resumeItem{Developed targeted corporate security awareness training materials based on data insights to mitigate future phishing risks.}
\resumeItem{\small\href{https://tinyurl.com/mastercardExp}{Forage certificate}}
\resumeItemListEnd
\resumeSubHeadingListEnd

%-----------TECHNICAL SKILLS-----------
\section{Technical Skills}
\begin{itemize}[leftmargin=0.15in, label={}]
\small
\item{
\textbf{Security}{: Security Monitoring, Threat Detection, Incident Response, Detection Engineering, Network Forensics, Malware Analysis} \\

\textbf{Security Tools}{: Splunk (SPL, saved searches), Elastic Stack (ELK), Zeek, Suricata, Wireshark/tshark, Sysmon, Winlogbeat, Wazuh, YARA, Sigma} \\

\textbf{Query Languages}{: Splunk SPL, KQL (Microsoft Sentinel), SQL} \\

\textbf{Operating Systems}{: Linux, Windows, Active Directory Administration} \\

\textbf{Programming}{: Python, Bash, SQL} \\

\textbf{Developer Tools}{: Git, GitHub, VS Code, PyCharm} \\

\textbf{Machine Learning}{: TensorFlow, Scikit-learn, Ensemble Learning (research)}
}
\end{itemize}
\end{document}
